%% GENERATED by tools/gen_latex.py -- edit content/ch06.py instead.
%% Build:  latexmk -pdf ch06.tex     (or: pdflatex ch06.tex, twice)
\documentclass[aspectratio=169,11pt,t]{beamer}
\usepackage{itbpro}

\coursetitle{Alur Kerja Universal Machine Learning}
\coursesubtitle{Anda tidak memulai dari kumpulan data. Anda memulai dari persoalan -- dan pekerjaan tersulitnya justru sebelum baris kode pertama.}
\coursekicker{Chapter 6}
\coursesource{Chollet \& Watson, Deep Learning with Python 3e -- bab 6 (hlm. 171-189)}
\courseduration{3 jam (2 sesi)}
\coursepresenter{Prof. Bambang Riyanto Trilaksono}{Pengajar Utama}
\coursebrand{AI for Professional \textperiodcentered\ ITB \textperiodcentered\ Ch 6}

\title{Alur Kerja Universal Machine Learning}
\author{Prof. Bambang Riyanto Trilaksono}
\date{}

\begin{document}
\covertitle

\framekicker{Learning outcomes}
\begin{frame}[shrink=25]{Tujuan Sesi}
\leadpar{Setelah sesi ini, peserta dapat:}
\begin{isteps}
\item Merumuskan sebuah persoalan bisnis menjadi \textbf{jenis tugas machine learning} yang tepat -- termasuk mengenali kapan machine learning \textbf{bukan} jawabannya.
\item Menyebut \textbf{dua hipotesis} yang diam-diam Anda buat setiap kali memulai proyek, dan apa artinya bila keduanya salah.
\item Merancang pengumpulan dan anotasi data yang \textbf{mewakili data produksi}, dan mengenali \textbf{bias pencuplikan}, \textbf{kebocoran target}, serta \textbf{pergeseran konsep}.
\item Memilih ukuran keberhasilan, lalu memilih \textbf{aktivasi lapis akhir, fungsi rugi, dan metrik} yang sepadan dengannya.
\item Menjalankan tiga tahap pengembangan model: \textbf{kalahkan tolok banding $\rightarrow$ besarkan sampai overfit $\rightarrow$ regularisasi dan setel}.
\item Memilih cara penyerahan -- \textbf{REST API, di perangkat, atau di peramban} -- beserta pengoptimalan inferensinya (pemangkasan dan kuantisasi bobot).
\item Menyusun \textbf{penetapan harapan pemangku kepentingan} dalam bahasa false-positive dan false-negative, bukan 'akurasi 98\%'.
\end{isteps}

\vskip 6pt
\vskip 4pt
\reslink{Buku}{Bab 6 --- teks penuh}{https://deeplearningwithpython.io/chapters/chapter06\_universal-workflow-of-ml}
\reslink{Notebook}{01\_memahami\_data\_sebelum\_model.ipynb}{../../notebooks/ch06/01\_memahami\_data\_sebelum\_model.ipynb}
\reslink{Notebook}{02\_pra\_pemrosesan\_umum.ipynb}{../../notebooks/ch06/02\_pra\_pemrosesan\_umum.ipynb}
\reslink{Notebook}{03\_ekspor\_dan\_kuantisasi.ipynb}{../../notebooks/ch06/03\_ekspor\_dan\_kuantisasi.ipynb}
\vskip 2pt
\end{frame}

\framekicker{Pembuka}
\begin{frame}[shrink=25]{Di dunia nyata, \inlinecode{keras.datasets} tidak ada}
\leadpar{Bayangkan Anda membuka jasa konsultasi machine learning sendiri. Proyeknya mulai berdatangan -- dan \hilite{tidak satu pun datang bersama kumpulan datanya}}
\begin{minipage}[t]{0.241\textwidth}
\begin{icard}
\cardh{Pencarian foto}
Ketik \textit{wedding}, dapatkan semua foto pernikahan -- tanpa penandaan manual.
\end{icard}
\end{minipage}%
\hfill
\begin{minipage}[t]{0.241\textwidth}
\begin{icard}
\cardh{Spam \& konten kasar}
Menandai kiriman pada aplikasi obrolan yang baru tumbuh.
\end{icard}
\end{minipage}%
\hfill
\begin{minipage}[t]{0.241\textwidth}
\begin{icard}
\cardh{Perekomendasi musik}
Untuk pengguna radio daring.
\end{icard}
\end{minipage}%
\hfill
\begin{minipage}[t]{0.241\textwidth}
\begin{icard}
\cardh{Fraud kartu kredit}
Untuk sebuah situs niaga elektronik.
\end{icard}
\end{minipage}%


\vskip 4pt

\begin{minipage}[t]{0.241\textwidth}
\begin{icard}
\cardh{Click-through rate iklan}
Memutuskan iklan mana yang disajikan ke siapa, kapan.
\end{icard}
\end{minipage}%
\hfill
\begin{minipage}[t]{0.241\textwidth}
\begin{icard}
\cardh{Kue cacat}
Menandai kue janggal di ban berjalan pabrik.
\end{icard}
\end{minipage}%
\hfill
\begin{minipage}[t]{0.241\textwidth}
\begin{icard}
\cardh{Situs arkeologi}
Menebak lokasi situs yang belum diketahui dari citra satelit.
\end{icard}
\end{minipage}%
\hfill
\begin{minipage}[t]{0.241\textwidth}
\begin{icardaccent}
\cardh{\ldots{}dan kasus Anda sendiri}
Yang akan dibawa ke tugas kelompok nanti.
\end{icardaccent}
\end{minipage}%
\begin{iband}[signal]
Ketujuh contoh ini dipakai berulang sepanjang bab, dan tiap tahap alur kerja diuji terhadap ketujuhnya. \hilite{Perhatikan bahwa dua di antaranya ternyata bukan persoalan deep learning sama sekali.}
\end{iband}
\end{frame}
\note{Minta peserta memilih satu kasus dari organisasinya sendiri di awal sesi lalu membawanya melewati setiap tahap di bab ini. Itu rangka tugas kelompoknya.}

\framekicker{Peta bab}
\begin{frame}[shrink=25]{Tiga bagian, dan yang tersulit ada di depan}
\begin{center}
\adjustbox{max width=\linewidth,max totalheight=0.52\textheight}{%

\begin{tikzpicture}[font=\sffamily\tiny,
  ar/.style={-{Stealth[length=4pt]}, signal, line width=0.8pt},
  bk/.style={-{Stealth[length=4pt]}, amberbr, line width=0.9pt, dashed}]
  \node[draw=itbbluelt!70, fill=itbbluelt!10, rounded corners=5pt, minimum width=3.3cm,
        minimum height=2.2cm, anchor=north west, align=left] (a) at (0,0)
    {\textbf{1 $\cdot$ Tetapkan tugas}\\[3pt]
     $\cdot$ rumuskan persoalannya\\$\cdot$ kumpulkan \& anotasi data\\
     $\cdot$ pahami datanya\\$\cdot$ pilih ukuran keberhasilan\\[2pt]
     \textcolor{amber}{bagian tersulit}};
  \node[draw=signal!70, fill=signal!10, rounded corners=5pt, minimum width=3.3cm,
        minimum height=2.2cm, anchor=north west, align=left] (b) at (3.7,0)
    {\textbf{2 $\cdot$ Kembangkan model}\\[3pt]
     $\cdot$ siapkan data\\$\cdot$ pilih protokol evaluasi\\
     $\cdot$ kalahkan tolok banding\\$\cdot$ besarkan sampai overfit\\
     $\cdot$ regularisasi \& setel};
  \node[draw=lime!70, fill=limebr!10, rounded corners=5pt, minimum width=3.3cm,
        minimum height=2.2cm, anchor=north west, align=left] (c) at (7.4,0)
    {\textbf{3 $\cdot$ Serahkan}\\[3pt]
     $\cdot$ jelaskan ke pemangku\\$\cdot$ kirim model inferensi\\
     $\cdot$ pantau di lapangan\\$\cdot$ rawat \& kumpulkan data baru};
  \draw[ar] (a.east) -- (b.west);
  \draw[ar] (b.east) -- (c.west);
  \draw[bk] (c.south) -- ++(0,-0.5) -| (a.south);
  \node[text=amber, anchor=north] at (5.4,-2.75)
    {data baru dari produksi --- menjadi bahan generasi model berikutnya};
\end{tikzpicture}

}
\end{center}
\vskip 3pt{\color{ink3}\fontsize{7.4}{9.4}\selectfont Alur kerja universal -- dan panah balik yang membuatnya tidak pernah benar-benar selesai.\par}
\begin{iquote}
\quotetext{Pengembangan model hanyalah satu langkah dalam alur kerja machine learning, dan menurut kami bukan yang tersulit. Yang paling sulit adalah \textbf{merumuskan persoalan serta mengumpulkan, menganotasi, dan membersihkan data}.}\quotecite{Chollet \& Watson, bab 6.2}
\end{iquote}
\end{frame}

\coursesection{01}{Menetapkan tugas}{Anda tidak bisa bekerja baik tanpa memahami konteksnya secara mendalam.}

\framekicker{Bagian 6.1.1}
\begin{frame}[shrink=25]{Empat pertanyaan yang harus ada di kepala}
\begin{isteps}
\item \textbf{Apa data masukannya? Apa yang hendak diramalkan?} Anda hanya bisa belajar meramalkan sesuatu bila ada data latihnya. \hilite{Ketersediaan data biasanya jadi faktor pembatas di tahap ini.}
\item \textbf{Jenis tugas machine learning apa ini?} Biner? Multikelas? Regresi skalar? Segmentasi citra? Perangkingan? Atau -- mungkin saja -- \textbf{machine learning bukan cara terbaik}, dan analisis statistik biasa lebih tepat.
\item \textbf{Seperti apa penyelesaian yang sudah ada?} Mungkin pelanggan sudah punya algoritma buatan tangan dengan banyak \inlinecode{if} bersarang. Mungkin ada manusia yang sekarang mengerjakannya secara manual. Pahami sistem yang sudah berjalan.
\item \textbf{Adakah kendala khusus?} Aplikasi terenkripsi ujung-ke-ujung $\rightarrow$ model harus hidup di ponsel pengguna. Kendala latensi ketat $\rightarrow$ model harus jalan di perangkat tertanam di pabrik, bukan di server jauh.
\end{isteps}
\begin{iband}[amber]
Perhatikan pertanyaan ketiga. Di banyak proyek nyata, \hilite{tolok bandingnya bukan 'acak', melainkan sistem berbasis aturan yang sudah berjalan bertahun-tahun} -- dan itu jauh lebih sulit dikalahkan.
\end{iband}
\end{frame}

\framekicker{Bagian 6.1.1}
\begin{frame}[shrink=25]{Ketujuh contoh, dipetakan ke jenis tugasnya}
\begin{itable}{>{\raggedright\arraybackslash}p{0.2256\textwidth}>{\raggedright\arraybackslash}p{0.282\textwidth}>{\raggedright\arraybackslash}p{0.4324\textwidth}}
\thead{Proyek} & \thead{Jenis tugasnya} & \thead{Catatan} \\ \midrule
Pencarian foto & Klasifikasi \textbf{multikelas, multilabel} & --- \\
Spam & Klasifikasi \textbf{biner} & Jadi \textbf{tiga kelas} kalau \textit{konten kasar} dijadikan kelas tersendiri. \\
Perekomendasi musik & \hilite{Bukan deep learning} & Lebih baik ditangani \textbf{faktorisasi matriks} (collaborative filtering). \\
Fraud kartu kredit & Klasifikasi \textbf{biner} & --- \\
Click-through rate & \textbf{Regresi skalar} & --- \\
Kue cacat & Klasifikasi \textbf{biner} & Tetapi butuh \textbf{deteksi objek} dulu untuk memotong kuenya dari citra mentah. Catatan buku: teknik yang dikenal sebagai \textit{anomaly detection} \hilite{justru tidak cocok di sini} \\
Situs arkeologi & \textbf{Perangkingan kemiripan citra} & Mengambil citra yang paling mirip situs yang sudah dikenal. \\
\end{itable}
\begin{iband}[signal]
Dua pelajaran dari tabel ini: ada persoalan yang \textbf{bukan} deep learning, dan ada persoalan yang \textbf{butuh dua model bertahap}. Keduanya sering terlewat kalau langsung meloncat ke pemodelan.
\end{iband}
\end{frame}

\framekicker{Bagian 6.1.1}
\begin{frame}[shrink=25]{Dua hipotesis yang selalu Anda buat diam-diam}
\begin{minipage}[t]{0.494\textwidth}
\begin{icardaccent}
\cardh{Target bisa diramalkan dari masukan}
Bahwa hubungan itu memang ada.
\end{icardaccent}
\end{minipage}%
\hfill
\begin{minipage}[t]{0.494\textwidth}
\begin{icardaccent}
\cardh{Data yang ada cukup informatif}
Untuk mempelajari hubungan antara masukan dan target itu.
\end{icardaccent}
\end{minipage}%
\begin{iband}[rose]
Sampai Anda punya model yang bekerja, keduanya \textbf{sekadar hipotesis} yang menunggu dibuktikan atau digugurkan. Menyusun contoh X dan target Y \hilite{tidak berarti X memuat cukup informasi untuk meramal Y}
\end{iband}
{\fontsize{9.4}{12.6}\selectfont\color{fg2}Contoh yang diberikan buku: meramalkan pergerakan saham dari riwayat harganya saja \textbf{tidak mungkin berhasil}, sebab riwayat harga tidak memuat banyak informasi yang meramalkan.\par}\vskip 4pt
{\fontsize{9.4}{12.6}\selectfont\color{fg2}Kalau setelah mencoba beberapa arsitektur yang masuk akal Anda tetap tidak bisa mengalahkan tolok banding sederhana, besar kemungkinan \textbf{jawaban atas pertanyaan Anda memang tidak ada di dalam data masukannya}. Kembali ke papan gambar.\par}\vskip 4pt
\end{frame}
\note{Ini slide yang menyelamatkan waktu paling banyak. Uji dua hipotesis ini di rapat perumusan, bukan setelah tiga bulan pemodelan.}

\framekicker{Bagian 6.1.1 $\cdot$ catatan etika}
\begin{frame}[shrink=25]{Teknologi tidak pernah netral}
{\fontsize{9.4}{12.6}\selectfont\color{fg2}Buku ini menyisipkan satu catatan etika, dan letaknya sengaja di tahap perumusan -- bukan di akhir. Contohnya: \textit{"membangun AI yang menilai tingkat dapat-dipercayanya seseorang dari foto wajahnya"}.\par}\vskip 4pt
\begin{isteps}
\item \textbf{Kesahihannya sendiri meragukan} -- tidak jelas mengapa sifat dapat-dipercaya akan tercermin di wajah seseorang.
\item Mengumpulkan datanya sama saja dengan \textbf{merekam bias dan prasangka orang-orang yang melabeli fotonya}.
\item Model yang dilatih di atasnya hanya akan \textbf{menyandikan bias yang sama ke dalam algoritma kotak hitam} -- yang justru memberinya lapisan tipis keabsahan.
\end{isteps}
\begin{iquote}
\quotetext{Di masyarakat yang sebagian besar belum melek teknologi seperti kita, \textit{"algoritma AI mengatakan orang ini tidak bisa dipercaya"} anehnya terdengar lebih berbobot dan lebih objektif daripada \textit{"John Smith mengatakan orang ini tidak bisa dipercaya"} -- padahal yang pertama adalah hampiran terpelajar atas yang kedua.}\quotecite{Chollet \& Watson, bab 6.1.1}
\end{iquote}
\begin{iband}[rose]
\textbf{Teknologi tidak pernah netral.} Kalau pekerjaan Anda berdampak pada dunia, dampak itu punya arah moral: \hilite{pilihan teknis juga pilihan etis}. Selalu sengaja memilih nilai apa yang hendak didukung pekerjaan Anda.
\end{iband}
\end{frame}
\note{Contoh yang mudah didekatkan ke peserta mana pun: penilaian kelayakan, penetapan harga, dan penandaan pelanggan --- semuanya bisa mencuci bias historis lalu memberinya wajah objektif.}

\framekicker{Bagian 6.1.2}
\begin{frame}[shrink=25]{Mengumpulkan data -- bagian termahal, dan yang paling berimbal hasil}
\begin{iquote}
\quotetext{Kalau Anda punya tambahan 50 jam untuk sebuah proyek, kemungkinan besar cara paling efektif membelanjakannya adalah \textbf{mengumpulkan lebih banyak data}, bukan mencari perbaikan pemodelan yang sedikit demi sedikit.}\quotecite{Chollet \& Watson, bab 6.1.2}
\end{iquote}
{\fontsize{9.4}{12.6}\selectfont\color{fg2}Titik bahwa \textbf{data lebih penting daripada algoritma} paling terkenal dikemukakan makalah peneliti Google tahun 2009, \textit{"The Unreasonable Effectiveness of Data"} -- judulnya memelesetkan buku Eugene Wigner 1960, \textit{The Unreasonable Effectiveness of Mathematics in the Natural Sciences}. Itu ditulis \textbf{sebelum} deep learning populer, dan naiknya deep learning \hilite{justru menambah} pentingnya data.\par}\vskip 4pt
\begin{itable}{>{\raggedright\arraybackslash}p{0.2444\textwidth}>{\raggedright\arraybackslash}p{0.3384\textwidth}>{\raggedright\arraybackslash}p{0.3572\textwidth}}
\thead{Pilihan anotasi} & \thead{Untungnya} & \thead{Risikonya} \\ \midrule
\textbf{Anotasi sendiri} & Kendali penuh atas mutu. & Lambat dan mahal dalam waktu. \\
\textbf{Urun daya} (mis. Mechanical Turk) & Murah dan berskala baik. & Anotasinya bisa \hilite{cukup berderau} \\
\textbf{Perusahaan pelabelan khusus} & Menghemat waktu dan biaya. & Melepaskan kendali. \\
\end{itable}
\begin{ibullets}
\item Apakah pelabelnya \textbf{harus ahli bidang}? Kucing vs anjing bisa siapa saja; ras anjing perlu pengetahuan khusus; CT scan patah tulang praktis menuntut gelar kedokteran.
\item Kalau perlu keahlian, \textbf{bisakah orang dilatih} untuk itu? Kalau tidak, bagaimana Anda mengakses ahlinya?
\item \textbf{Apakah Anda sendiri paham} bagaimana ahli itu sampai pada anotasinya? Kalau tidak, data Anda jadi kotak hitam dan rekayasa fitur manual tertutup -- tidak fatal, tetapi membatasi.
\item Perangkat lunak anotasi yang produktif \textbf{menghemat sangat banyak waktu}; layak diinvestasikan sejak awal proyek.
\end{ibullets}
\end{frame}

\framekicker{Bagian 6.1.2}
\begin{frame}[shrink=25]{Data yang tidak mewakili -- dosa besar}
\begin{minipage}[t]{0.485\textwidth}
{\fontsize{9.4}{12.6}\selectfont\color{fg2}\textbf{Contoh dari buku.} Aplikasi pengenal masakan dari foto. Model dilatih dengan foto dari jejaring sosial berbagi gambar yang populer di kalangan penggemar kuliner. Akurasi uji jauh di atas 90\%.\par}\vskip 4pt
\begin{iband}[rose]
Setelah dirilis: \textbf{salah 8 dari 10 kali}. Foto unggahan pengguna -- masakan acak, restoran acak, ponsel acak -- \hilite{sama sekali tidak menyerupai} foto profesional yang pencahayaannya bagus dan menggugah selera itu.
\end{iband}

\end{minipage}%
\hfill
\begin{minipage}[t]{0.485\textwidth}
{\fontsize{9.4}{12.6}\selectfont\color{fg2}\textbf{Aturannya}\par}\vskip 4pt
\begin{ibullets}
\item Bila mungkin, kumpulkan data \textbf{langsung dari lingkungan tempat model akan dipakai}.
\item Model sentimen ulasan film dipakai pada ulasan IMDB baru -- bukan pada ulasan restoran Yelp, bukan pada status Twitter.
\item Kalau melatih di data produksi tidak mungkin, \textbf{pahami betul bedanya} dan koreksi perbedaan itu secara aktif.
\end{ibullets}

\end{minipage}%
{\fontsize{9.4}{12.6}\selectfont\color{fg2}\textbf{Bias pencuplikan} adalah bentuknya yang paling licik: proses pengumpulan data berinteraksi dengan hal yang hendak Anda ramalkan. Contoh sejarahnya terkenal -- malam pemilu 1948, \textit{Chicago Tribune} memasang tajuk \textbf{"DEWEY DEFEATS TRUMAN"}. Paginya Truman yang menang. Redakturnya memercayai survei telepon; padahal pengguna telepon pada 1948 \hilite{bukan cuplikan acak yang mewakili pemilih} -- mereka cenderung lebih kaya dan konservatif.\par}\vskip 4pt
\end{frame}

\framekicker{Bagian 6.1.2}
\begin{frame}[shrink=25]{Pergeseran konsep -- dan umur pakai model}
\begin{center}
\adjustbox{max width=\linewidth,max totalheight=0.52\textheight}{%

\begin{tikzpicture}[font=\sffamily\tiny]
  \node[text=ink2, anchor=east] at (0,1.4) {deteksi fraud kartu};
  \node[draw=rose!75, fill=rosebr!22, rounded corners=2.5pt, minimum width=0.8cm,
        minimum height=0.34cm, anchor=west] at (0.15,1.4) {};
  \node[text=rose, anchor=west, font=\ttfamily\tiny] at (1.1,1.4) {hari};

  \node[text=ink2, anchor=east] at (0,0.8) {perekomendasi musik};
  \node[draw=amber!70, fill=amberbr!22, rounded corners=2.5pt, minimum width=1.9cm,
        minimum height=0.34cm, anchor=west] at (0.15,0.8) {};
  \node[text=amber, anchor=west, font=\ttfamily\tiny] at (2.2,0.8) {minggu};

  \node[text=ink2, anchor=east] at (0,0.2) {mesin pencari citra};
  \node[draw=lime!65, fill=limebr!20, rounded corners=2.5pt, minimum width=5.6cm,
        minimum height=0.34cm, anchor=west] at (0.15,0.2) {};
  \node[text=lime, anchor=west, font=\ttfamily\tiny] at (5.9,0.2) {beberapa tahun (paling baik)};

  \draw[rule] (0.15,-0.2) -- (8.6,-0.2);
  \node[text=ink3, anchor=north] at (4.4,-0.3) {umur pakai sebelum model harus dilatih ulang};
  \node[text=amber, anchor=west, align=left] at (-2.4,-1.0)
    {Pergeseran konsep paling tajam di konteks adversarial:\\pola kecurangan berubah nyaris tiap hari.};
\end{tikzpicture}

}
\end{center}
\vskip 3pt{\color{ink3}\fontsize{7.4}{9.4}\selectfont Angka-angka ini dari buku. Perhatikan bahwa deteksi fraud -- kasus yang paling adversarial -- adalah yang paling cepat basi.\par}
\begin{ibullets}
\item \textbf{Pergeseran konsep} terjadi saat sifat data produksi berubah seiring waktu, sehingga akurasi model merosot perlahan.
\item Kumpulan IMDB dikumpulkan pada 2011; model yang dilatih di atasnya akan bekerja lebih buruk pada ulasan 2020 daripada ulasan 2012 -- kosakata, ungkapan, dan genre film berubah.
\item Menangani pergeseran yang cepat menuntut \textbf{pengumpulan data, anotasi, dan pelatihan ulang yang terus-menerus}.
\end{ibullets}
\begin{iband}[amber]
Kalimat penutupnya keras: memakai machine learning yang dilatih atas data masa lalu untuk meramal masa depan berarti \textbf{mengandaikan masa depan akan berkelakuan seperti masa lalu}. \hilite{Sering kali tidak.}
\end{iband}
\end{frame}

\framekicker{Bagian 6.1.3}
\begin{frame}[shrink=25]{Memahami data Anda -- daftar periksa sebelum melatih}
{\fontsize{9.4}{12.6}\selectfont\color{fg2}Memperlakukan kumpulan data sebagai kotak hitam adalah \textbf{praktik yang buruk}. Sebelum melatih apa pun, jelajahi dan gambarkan datanya.\par}\vskip 4pt
\begin{isteps}
\item Ada citra atau teks? \textbf{Lihat langsung} beberapa contohnya, berikut labelnya.
\item Ada fitur numerik? \textbf{Gambar histogramnya} untuk merasakan rentang nilai dan seberapa sering tiap nilai muncul.
\item Ada informasi lokasi? \textbf{Petakan}. Apakah ada pola yang muncul?
\item Ada sampel yang \textbf{nilainya hilang}? Itu harus ditangani saat penyiapan.
\item Tugas klasifikasi? \textbf{Cetak jumlah contoh tiap kelas.} Kalau tidak seimbang, ketidakseimbangan itu harus diperhitungkan.
\item Periksa \textbf{kebocoran target}.
\end{isteps}
\begin{iband}[rose]
\textbf{Kebocoran target} -- adanya fitur yang memberi informasi tentang target tetapi \hilite{tidak akan tersedia di produksi}. Contoh buku: melatih model atas rekam medis untuk meramal apakah seseorang akan dirawat karena kanker, sementara rekamnya memuat fitur \textit{"orang ini telah didiagnosis kanker"}.
\end{iband}
{\fontsize{9.4}{12.6}\selectfont\color{fg2}Pertanyaan yang harus selalu Anda ajukan: \textbf{apakah setiap fitur dalam data saya akan tersedia dalam bentuk yang sama di produksi?}\par}\vskip 4pt
\end{frame}
\note{Bentuk yang paling sering: kolom yang diisi BELAKANGAN oleh petugas --- status penanganan, kode penutupan, catatan tindak lanjut. Semuanya belum ada saat model harus memutuskan.}

\framekicker{Bagian 6.1.4}
\begin{frame}[shrink=25]{Memilih ukuran keberhasilan}
\begin{iquote}
\quotetext{Untuk mengendalikan sesuatu, Anda harus bisa mengamatinya. Untuk berhasil dalam sebuah proyek, Anda harus lebih dulu menetapkan apa yang Anda maksud dengan berhasil.}\quotecite{Chollet \& Watson, bab 6.1.4}
\end{iquote}
\begin{itable}{>{\raggedright\arraybackslash}p{0.376\textwidth}>{\raggedright\arraybackslash}p{0.564\textwidth}}
\thead{Bentuk persoalan} & \thead{Metrik yang lazim} \\ \midrule
Klasifikasi \textbf{seimbang} & Akurasi, dan \textbf{AUC} dari kurva ROC. \\
Kelas \textbf{tidak seimbang}, perangkingan, multilabel & \textbf{Presisi dan recall}, atau metrik yang menghitung false positive / true positive / false negative / true negative. \\
Bukan salah satu di atas & Tidak jarang Anda harus \textbf{menetapkan metrik sendiri}. \\
\end{itable}
\begin{iband}[signal]
Metrik keberhasilan \textbf{menuntun semua pilihan teknis} sepanjang proyek. Ia harus selaras langsung dengan sasaran yang lebih tinggi -- yaitu \hilite{keberhasilan bisnis pelanggan Anda}.
\end{iband}
{\fontsize{9.4}{12.6}\selectfont\color{fg2}Untuk merasakan keragaman metrik dan kaitannya dengan ranah persoalan, buku menyarankan menjelajahi kompetisi di \hlink{https://kaggle.com}{Kaggle}.\par}\vskip 4pt
\end{frame}

\coursesection{02}{Mengembangkan model}{Bagian yang paling banyak diajarkan tutorial -- dan bukan yang tersulit.}

\framekicker{Bagian 6.2.1}
\begin{frame}[shrink=25]{Menyiapkan data: vektorisasi, normalisasi, nilai hilang}
\begin{minipage}[t]{0.485\textwidth}
{\fontsize{9.4}{12.6}\selectfont\color{fg2}\textbf{Vektorisasi.} Semua masukan dan target harus berupa tensor bilangan titik-mengambang (atau, dalam kasus khusus, tensor bilangan bulat atau string). Suara, citra, teks -- semuanya diubah jadi tensor lebih dulu.\par}\vskip 4pt
{\fontsize{9.4}{12.6}\selectfont\color{fg2}\textbf{Normalisasi nilai.} Tidak aman menyuapkan data bernilai relatif besar (bilangan bulat berdigit banyak, jauh lebih besar dari nilai awal bobot jaringan) atau data yang \textbf{heterogen} (satu fitur di 0-1, yang lain di 100-200). Itu memicu pembaruan gradien besar yang \hilite{mencegah jaringan konvergen}\par}\vskip 4pt

\end{minipage}%
\hfill
\begin{minipage}[t]{0.485\textwidth}
\icode{python}{dua sifat yang diinginkan}{listings/ch06-01.py}

\end{minipage}%
\begin{itable}{>{\raggedright\arraybackslash}p{0.2444\textwidth}>{\raggedright\arraybackslash}p{0.6956\textwidth}}
\thead{Nilai hilang pada fitur\ldots{}} & \thead{Yang dilakukan} \\ \midrule
\textbf{Kategorik} & Aman membuat \textbf{kategori baru} yang berarti \textit{nilainya hilang}. Model akan belajar sendiri apa artinya terhadap target. \\
\textbf{Numerik} & \textbf{Hindari mengisi nilai sembarang seperti 0} -- itu bisa menciptakan ketakbersinambungan pada ruang laten. Pakai \textbf{rerata atau median} fitur itu; atau latih model untuk menebak nilainya dari fitur lain. \\
\end{itable}
\begin{iband}[amber]
Jebakan yang halus: kalau Anda \textbf{mengharapkan} fitur kategorik hilang di data uji tetapi jaringan dilatih tanpa satu pun nilai hilang, jaringan \hilite{tidak pernah belajar mengabaikannya}. Obatnya: buat sampel latih buatan yang bolong -- salin beberapa sampel, lalu buang fitur yang diperkirakan akan hilang.
\end{iband}
\end{frame}

\framekicker{Bagian 6.2.3 $\cdot$ tabel 6.1}
\begin{frame}[shrink=25]{Kalahkan tolok banding: tiga hal yang perlu benar}
\begin{minipage}[t]{0.3253\textwidth}
\begin{icardaccent}
\cardh{Rekayasa fitur}
Saring fitur yang tidak informatif (seleksi fitur), dan pakai pengetahuan Anda tentang persoalan untuk membuat fitur baru.
\end{icardaccent}
\end{minipage}%
\hfill
\begin{minipage}[t]{0.3253\textwidth}
\begin{icardaccent}
\cardh{Prior arsitektur yang benar}
Jaringan padat? ConvNet? Rekuren? Transformer? Atau -- apakah deep learning memang pendekatan yang baik untuk tugas ini?
\end{icardaccent}
\end{minipage}%
\hfill
\begin{minipage}[t]{0.3253\textwidth}
\begin{icardaccent}
\cardh{Konfigurasi pelatihan yang cukup baik}
Fungsi rugi apa? Ukuran batch dan learning rate berapa?
\end{icardaccent}
\end{minipage}%
\begin{itable}{>{\raggedright\arraybackslash}p{0.2444\textwidth}>{\raggedright\arraybackslash}p{0.1692\textwidth}>{\raggedright\arraybackslash}p{0.2444\textwidth}>{\raggedright\arraybackslash}p{0.282\textwidth}}
\thead{Tugas} & \thead{Aktivasi lapis akhir} & \thead{Fungsi rugi} & \thead{Metrik} \\ \midrule
Klasifikasi biner & Sigmoid & Binary crossentropy & Binary accuracy, ROC AUC \\
Multikelas, label tunggal & Softmax & Categorical crossentropy & Categorical accuracy, top-k, ROC AUC \\
Multikelas, multilabel & Sigmoid & Binary crossentropy & Binary accuracy, ROC AUC \\
Regresi & Tidak ada & Mean squared error & Mean absolute error \\
\end{itable}
\begin{iband}[signal]
\textbf{Mengapa metriknya tidak langsung dioptimalkan?} Fungsi rugi harus bisa dihitung hanya dari satu mini-batch (idealnya bahkan dari satu titik data) dan harus \textbf{terdiferensialkan}. ROC AUC tidak memenuhi keduanya, jadi yang dioptimalkan adalah \hilite{pengganti}-nya, biasanya crossentropy -- dengan harapan makin rendah crossentropy, makin tinggi ROC AUC.
\end{iband}
\end{frame}

\framekicker{Bagian 6.2.4-6.2.5}
\begin{frame}[shrink=25]{Besarkan, lalu regularisasi}
\begin{minipage}[t]{0.485\textwidth}
{\fontsize{9.4}{12.6}\selectfont\color{fg2}\textbf{6.2.4 $\cdot$ Kembangkan model yang overfit}\par}\vskip 4pt
\begin{ibullets}
\item Tambah lapis.
\item Perbesar lapisnya.
\item Latih lebih banyak epoch.
\end{ibullets}
\begin{iband}[signal]
Model ideal berdiri \textbf{tepat di perbatasan} antara underfit dan overfit. Untuk tahu di mana perbatasannya, \hilite{Anda harus melewatinya lebih dulu.}
\end{iband}

\end{minipage}%
\hfill
\begin{minipage}[t]{0.485\textwidth}
{\fontsize{9.4}{12.6}\selectfont\color{fg2}\textbf{6.2.5 $\cdot$ Regularisasi dan setel} -- tahap yang makan waktu paling banyak\par}\vskip 4pt
\begin{ibullets}
\item Coba arsitektur berbeda; tambah atau kurangi lapis.
\item Tambahkan \textbf{dropout}.
\item Kalau modelnya kecil, tambahkan \textbf{regularisasi L1 atau L2}.
\item Coba hiperparameter lain -- jumlah unit per lapis, learning rate.
\item Bila perlu, ulangi \textbf{kurasi data atau rekayasa fitur}.
\end{ibullets}

\end{minipage}%
\begin{iband}[rose]
Peringatan yang sama, sekali lagi: tiap kali Anda memakai umpan balik validasi untuk menyetel model, informasi bocor ke dalam model. Beberapa kali tidak apa-apa; \textbf{dilakukan sistematis selama banyak iterasi, model akan overfit terhadap proses validasi itu sendiri} -- walau tidak ada model yang dilatih langsung atas data validasi.
\end{iband}
{\fontsize{9.4}{12.6}\selectfont\color{fg2}Sebagian besar pekerjaan ini bisa diotomasi dengan perangkat penyetelan hiperparameter seperti \textbf{KerasTuner} (bab 18). Dan bila unjuk kerja di himpunan uji ternyata jauh lebih buruk daripada di validasi: prosedur validasi Anda tidak andal, atau Anda sudah overfit ke data validasi. \hilite{Pindah ke protokol yang lebih andal, misalnya K-lipat berulang.}\par}\vskip 4pt
\end{frame}

\coursesection{03}{Menyerahkan model}{Proyek tidak berakhir di notebook Colab yang bisa menyimpan model terlatih.}

\framekicker{Bagian 6.3.1}
\begin{frame}[shrink=25]{Menetapkan harapan -- jangan bilang 'akurasi 98\%'}
\begin{ibullets}
\item Harapan orang awam terhadap sistem AI sering \textbf{tidak realistis}: mereka mengira sistem \textit{memahami} tugasnya dan punya akal sehat seperti manusia.
\item Obatnya: \textbf{tunjukkan contoh cara model itu gagal} -- terutama kesalahan klasifikasi yang terasa mengejutkan.
\item Hindari pernyataan abstrak seperti \textit{"modelnya berakurasi 98\%"}, yang \hilite{oleh kebanyakan orang dibulatkan ke 100\%} dalam kepala mereka.
\end{ibullets}
\begin{iband}[signal]
Bicaralah dalam \textbf{laju false-negative dan false-positive}, lalu terjemahkan ke angka harian yang bisa dibayangkan.
\end{iband}
\begin{iquote}
\quotetext{Dengan setelan ini, model deteksi fraud akan punya laju false-negative 5\% dan false-positive 2,5\%. Setiap hari, rata-rata 200 transaksi sah akan ditandai sebagai fraud dan dikirim ke pemeriksaan manual, dan rata-rata 14 transaksi fraud akan terlewat. Rata-rata 266 transaksi fraud akan tertangkap dengan benar.}\quotecite{Contoh penetapan harapan dari Chollet \& Watson, bab 6.3.1}
\end{iquote}
\begin{iband}[amber]
Bahas juga \textbf{pemilihan ambang} bersama pemangku kepentingan -- ambang peluang yang berbeda menghasilkan laju false-negative dan false-positive yang berbeda. Keputusan itu menyangkut pertukaran yang \hilite{hanya bisa ditangani dengan pemahaman mendalam atas konteks bisnisnya}
\end{iband}
\end{frame}
\note{Ini slide yang paling langsung bisa dipakai peserta besok pagi. Latihan kelas: tulis ulang satu klaim akurasi dari proyek Anda ke dalam bentuk kalimat di atas.}

\framekicker{Bagian 6.3.2}
\begin{frame}[shrink=25]{Tiga cara menyerahkan model}
\begin{minipage}[t]{0.3253\textwidth}
\begin{icardaccent}
\cardh{REST API}
Pakai bila: ada \textbf{akses internet andal}; \textbf{tidak ada kendala latensi ketat} (pulang-pergi sekitar \textbf{500 ms}); dan data masukannya \textbf{tidak sangat peka} -- sebab data harus tersedia dalam bentuk terdekripsi di server.
\cardtag{pencari citra $\cdot$ perekomendasi $\cdot$ fraud $\cdot$ satelit}
\end{icardaccent}
\end{minipage}%
\hfill
\begin{minipage}[t]{0.3253\textwidth}
\begin{icardaccent}
\cardh{Di perangkat}
Pakai bila: \textbf{latensi ketat} atau \textbf{konektivitas rendah}; model bisa dibuat cukup kecil; akurasi tertinggi \textbf{bukan} hal kritis; dan data masukannya sangat peka sehingga tak boleh didekripsi di server jauh.
\cardtag{spam terenkripsi $\cdot$ kue di pabrik}
\end{icardaccent}
\end{minipage}%
\hfill
\begin{minipage}[t]{0.3253\textwidth}
\begin{icardaccent}
\cardh{Di peramban}
Pakai bila: ingin \textbf{memindahkan komputasi ke pengguna} (biaya server turun tajam); data harus tetap di perangkat pengguna; latensi ketat (hemat \textasciitilde{}100 ms pulang-pergi jaringan); atau aplikasi harus tetap jalan \textbf{tanpa koneksi}.
\cardtag{versi web \& desktop aplikasi obrolan}
\end{icardaccent}
\end{minipage}%
\begin{iband}[rose]
Peringatan untuk penyerahan di peramban: seluruh model \textbf{diunduh ke perangkat pengguna}. Pastikan tidak ada bagian model yang harus dirahasiakan -- sebab dari model terlatih \hilite{biasanya masih mungkin memulihkan sebagian informasi tentang data pelatihannya}. Jangan publikasikan model yang dilatih atas data peka.
\end{iband}
\end{frame}

\framekicker{Bagian 6.3.2}
\begin{frame}[shrink=25]{Mengekspor model: TensorFlow Serving dan ONNX}
\begin{minipage}[t]{0.485\textwidth}
\icode{python}{TensorFlow SavedModel}{listings/ch06-02.py}

\end{minipage}%
\hfill
\begin{minipage}[t]{0.485\textwidth}
\icode{python}{ONNX}{listings/ch06-03.py}

\end{minipage}%
{\fontsize{9.4}{12.6}\selectfont\color{fg2}Keduanya bekerja dengan \textbf{mengangkat seluruh bobot model dan graf komputasinya keluar dari program Python}, sehingga bisa dilayani dari banyak lingkungan berbeda -- misalnya server C++. Kalau ini terdengar mirip mekanisme kompilasi di bab 3, \hilite{memang begitu}: TensorFlow Serving pada dasarnya pustaka untuk melayani graf \inlinecode{tf.function} dengan sehimpunan bobot tersimpan.\par}\vskip 4pt
\begin{ibullets}
\item \textbf{TensorFlow Lite} -- inferensi di perangkat: Android, iOS, CPU ARM, Raspberry Pi, dan sebagian mikrokontroler. Formatnya sama dengan TensorFlow Serving. Runtime ONNX juga bisa jalan di perangkat bergerak.
\item \textbf{TensorFlow.js} -- menjalankan model di peramban; ia bahkan mengimplementasikan hampir seluruh API Keras (nama kerjanya dulu \textit{WebKeras}). ONNX punya runtime JavaScript sendiri.
\item Pilihan lain: \textbf{layanan awan terkelola} seperti Cloud AI Platform, yang mengurus pem-batch-an prediksi, penyeimbangan beban, dan penskalaan.
\end{ibullets}
\end{frame}

\framekicker{Bagian 6.3.2}
\begin{frame}[shrink=25]{Mengoptimalkan model untuk inferensi}
\begin{minipage}[t]{0.494\textwidth}
\begin{icardaccent}
\cardh{Pemangkasan bobot (weight pruning)}
Tidak setiap koefisien pada tensor bobot menyumbang sama besar terhadap prediksi. Dengan hanya menyimpan yang paling berarti, jumlah parameter bisa \textbf{turun banyak} dengan ongkos kecil pada metrik. Seberapa banyak dipangkas = kendali Anda atas pertukaran ukuran lawan akurasi.
\end{icardaccent}
\end{minipage}%
\hfill
\begin{minipage}[t]{0.494\textwidth}
\begin{icardaccent}
\cardh{Kuantisasi bobot (weight quantization)}
Model dilatih dengan bobot float32. Bobot itu bisa dikuantisasi ke \textbf{int8}, menghasilkan model khusus-inferensi yang \textbf{empat kali lebih kecil} tetapi tetap mendekati akurasi aslinya.
\end{icardaccent}
\end{minipage}%
\icode{python}{API kuantisasi bawaan Keras}{listings/ch06-04.py}
\begin{iband}[signal]
Pengoptimalan ini \textbf{terutama penting} saat menyerahkan ke lingkungan dengan kendala daya dan memori ketat -- ponsel dan perangkat tertanam -- atau untuk aplikasi berlatensi rendah. Lakukan \hilite{sebelum} mengimpor ke TensorFlow.js atau mengekspor ke TensorFlow Lite.
\end{iband}
\end{frame}

\framekicker{Bagian 6.3.3-6.3.4}
\begin{frame}[shrink=25]{Memantau dan merawat -- pekerjaan yang tak berujung}
{\fontsize{9.4}{12.6}\selectfont\color{fg2}Anda sudah mengekspor model inferensi, memasangkannya ke aplikasi, dan menjalankan uji coba di data produksi. \textbf{Bahkan ini belum akhirnya.}\par}\vskip 4pt
\begin{isteps}
\item \textbf{Uji A/B teracak.} Sebagian kasus lewat model baru, sebagian lagi -- kelompok kendali -- tetap lewat proses lama. Setelah cukup banyak kasus, selisih hasilnya \hilite{bisa dikaitkan ke modelnya}, bukan ke perubahan lain.
\item \textbf{Audit manual berkala} atas prediksi di data produksi. Infrastruktur anotasi yang sama bisa dipakai ulang: kirim sebagian data produksi untuk dianotasi manual, lalu bandingkan.
\item \textbf{Kalau audit manual mustahil}, cari jalur penilaian lain -- misalnya survei pengguna, untuk sistem penandaan spam dan konten kasar.
\end{isteps}
\begin{iband}[amber]
\textbf{Begitu model diluncurkan, Anda harus sudah bersiap melatih generasi berikutnya yang akan menggantikannya.}
\end{iband}
\begin{ibullets}
\item Awasi \textbf{perubahan pada data produksi}. Apakah ada fitur baru? Apakah himpunan labelnya perlu diperluas atau diubah?
\item Terus kumpulkan dan anotasi data, dan \textbf{perbaiki jalur anotasi} dari waktu ke waktu.
\item Beri perhatian khusus pada \textbf{sampel yang tampak sulit diklasifikasikan} oleh model sekarang -- sampel itulah yang \hilite{paling mungkin memperbaiki unjuk kerja}.
\end{ibullets}
\end{frame}

\framekicker{Ringkasan}
\begin{frame}[shrink=25]{Yang wajib terbawa dari bab 6}
\begin{isteps}
\item \textbf{Tetapkan tugasnya dulu.} Pahami konteks, tujuan akhir, dan kendalanya; kumpulkan dan anotasi data; pilih cara mengukur keberhasilan.
\item Anda selalu membuat \textbf{dua hipotesis}. Sampai ada model yang bekerja, keduanya belum terbukti.
\item \textbf{Pilihan teknis juga pilihan etis.} Teknologi tidak pernah netral.
\item \textbf{Data harus mewakili produksi.} Waspadai bias pencuplikan, kebocoran target, dan pergeseran konsep.
\item \textbf{Kalahkan tolok banding $\rightarrow$ besarkan sampai overfit $\rightarrow$ regularisasi dan setel.} Dalam urutan itu.
\item \textbf{Serahkan sesuai kendalanya}: REST API, di perangkat, atau di peramban -- lalu pangkas dan kuantisasi untuk inferensi.
\item \textbf{Bicara dalam false-positive dan false-negative}, bukan persen akurasi.
\item \textbf{Tidak ada model yang bertahan selamanya.} Pantau, audit, dan siapkan generasi berikutnya sejak hari peluncuran.
\end{isteps}
\vskip 4pt
\reslink{NOTEBOOK}{03\_ekspor\_dan\_kuantisasi.ipynb}{../../course-slides/notebooks/ch06/03\_ekspor\_dan\_kuantisasi.ipynb}
\reslink{BAB BERIKUT}{Bab 7 --- Menyelam ke Keras}{../ch07/index.html}
\vskip 2pt
\end{frame}

\end{document}
